\begin{center}
\section*{P}
\end{center}
\addcontentsline{toc}{section}{P}
\term{Piano di Progetto}
Documento che si occupa di definire l'organizzazione del gruppo, andando a specificare l'assegnazione dei ruoli, il monte ore previsto per ogni risorsa, le tempistiche e il costo economico preventivo in funzione dell'avanzamento del progetto.
\term{Piano di Qualifica}
Documento che si occupa di garantire il rispetto dei requisiti del progetto definendo le strategie di verifica e validazione affinchè vengano rispettate le attese del proponente.
\term{Product Baseline}
Fase finale del progetto che prevede la realizzazione di un MVP, a seguito del superamento dei test concordati nella Requirements and Technology Baseline.
\term{Progettazione}
Processo che definisce la struttura di un sistema, identificandone componenti, relazioni e livello di dettaglio necessario per la realizzazione.
\term{Progettista}
Ruolo dedito alla progettazione dell'architettura di un progetto, in modo da assicurarsi codice corretto e manutenibile.
\term{Progetto}
Insieme ordinato di attività, istanziate dai processi di ciclo di vita che soddisfano gli obiettivi dati.
\term{Programmatore}
Persona incaricata di svolgere le attività di codifica.
\term{Proof of Concept}
Prodotto dimostrativo atto a dimostrare la fattibilità dell'idea.
\term{Proponente}
Persona, gruppo o organizzazione che propone un progetto, idea o iniziativa e ne definisce obiettivi, requisiti e motivazioni.
